\section{Introduction}
Organic photovoltaics (OPVs) have attracted significant scientific attention as a low-cost, lightweight, and flexible technology for solar energy conversion. Over the past decade, power conversion efficiencies (PCEs) of OPVs have surpassed 18\%, driven primarily by the development of novel non-fullerene acceptors (NFAs) and high-performance donor polymers.\cite{Vance2025,Adams2024} However, to compete with silicon-based technologies, a deeper understanding of the fundamental photophysical processes governing charge generation, transport, and recombination in these systems is required.

In this work, we focus on donor-acceptor molecular architectures incorporating benzothiadiazole (BTD) units, which are known for their strong electron-withdrawing capabilities and broad absorption profiles in the visible spectrum. The donor molecules investigated here, Nexa-101 and Nexa-102 (supplied by Nexa Materials), feature a central BTD core coupled with terminal thiophene derivatives to facilitate intermolecular $\pi$-$\pi$ stacking and enhance charge carrier mobility. 

The primary objective of this study is to elucidate the excited-state dynamics of these BTD-based donors using femtosecond transient absorption (TA) spectroscopy. We investigate the pathways of singlet exciton decay, focusing on the potential for singlet fission to generate multiple charge carriers from a single absorbed photon. Through a collaborative effort with the Synapse Research Institute, we compare these experimental findings with quantum chemical simulations to establish a comprehensive structure-property relationship.
